\chapter{The Wage Structure}\label{ch:b3:wage-structure}

Wages differ between people with the same skills, and between people with
different ones. Some differences pay for unpleasant or risky work; some reward
skills that are scarce; some reflect who the worker is rather than what she can do;
some come from bargaining power. This chapter takes the four in turn and asks of
each how it is measured: the market for skills over a century, discrimination by
experiment and by decomposition, the union premium, and the effect of immigrants
on the wages of those already there.

\begin{openingchart}{The race between education and technology, United States,
1914--2017: annual changes in the college wage premium, in the relative supply of
college-educated workers and in their relative demand, by period (100 $\times$ annual log
change; demand computed with an elasticity of substitution of 1.62).}{Autor, Goldin
and Katz (2020), NBER Working Paper 26705, Table A2}
\begin{tikzpicture}
\begin{axis}[omrecord, ybar, bar width=7pt, width=0.74\linewidth,
  height=0.42\linewidth, xmin=-0.6, xmax=4.6, ymin=-1, ymax=4.5,
  xtick={0,...,4}, xticklabels from table={figdata/bachelor-3/25-wage-structure/out/race.csv}{period},
  x tick label style={font=\scriptsize}, ylabel={100 $\times$ log change a year},
  legend pos=north west, legend style={font=\footnotesize}]
  \addplot[fill=omThm, draw=omThm, area legend] table[x=pos, y=wage] {figdata/bachelor-3/25-wage-structure/out/race.csv};
  \addplot[fill=omDef!60, draw=omDef, area legend] table[x=pos, y=supply] {figdata/bachelor-3/25-wage-structure/out/race.csv};
  \addplot[fill=black!35, draw=black, area legend] table[x=pos, y=demand] {figdata/bachelor-3/25-wage-structure/out/race.csv};
  \legend{Premium, Relative supply, Relative demand}
\end{axis}
\end{tikzpicture}
\end{openingchart}

% ledger: b3ws:agk.a2
The supply of college-educated workers relative to others grew by 2.57~\% a year from 1914 to
1939 and by 3.51~\% from 1959 to 1979, and the premium fell or stood still. From 1979 to 1999
supply growth slowed to 2.28~\% while demand grew 4.21~\% a year, and the premium rose by 1.19
log points a year. After 1999 demand growth slowed too. Over the whole century demand
and supply grew at almost the same pace, and the premium ended near where it
began. A model with two kinds of workers and one parameter accounts for most of this.

\begin{recall}
The grades 10--12 volume defined the compensating differential and decomposed a
wage gap by hand; the Lorenz curve and Gini coefficient are also its. The
elasticity of substitution is the Year 1 volume's, and the human-capital model
\cref{ch:b3:human-capital}'s. Instruments and difference in differences are
chapters 1 and 3.
\end{recall}

\begin{notation}[Symbols of this chapter]\label{not:b3:wage-structure:symbols}
$S$ and $U$ are the numbers of skilled and unskilled workers, $w_s$ and $w_u$ their wages; $\sigma$
the elasticity of substitution between them; $D_t$ the log relative demand for skill;
$\bar x_g$ and $\beta_g$ group $g$'s mean characteristics and coefficients.
\end{notation}

\section{Compensating differences}\label{sec:b3:wage-structure:compensating}

\begin{definition}[Compensating wage differential]\label{def:b3:wage-structure:compensating}
\index{compensating wage differential}\index{hedonic wage}
A \emph{compensating wage differential} is the part of a wage that pays for a job's
non-wage features --- risk, effort, hours, place --- and makes the marginal worker
indifferent between jobs; across jobs the wage traces a hedonic locus, as house
prices did in \cref{ch:b3:cost-benefit-analysis}.
\end{definition}

A job with an extra fatality risk of 1 in 10{,}000 a year that pays 600 dalens more than a
safe one, to the marginal worker, implies a \omterm{def:b3:cost-benefit-analysis:vsl}{value of a statistical life} of 6 million
dalens: the \omterm{def:b3:cost-benefit-analysis:vsl}{value of a
statistical life} of \cref{ch:b3:cost-benefit-analysis} is estimated in exactly this way.
The danger is that the jobs differ in other ways too, and that workers who accept
risk are those who mind it least.

\section{The skill premium}\label{sec:b3:wage-structure:premium}

\begin{definition}[Skill premium; skill-biased technical change]\label{def:b3:wage-structure:skill-premium}
\index{skill premium}\index{college wage premium}\index{skill-biased technical change}
The \emph{skill premium} is the ratio of the wage of skilled to unskilled workers,
usually measured as the log wage gap between college graduates and those who
finished only high school (the \emph{college wage premium}). \emph{Skill-biased
technical change} is technical change that raises the relative demand for
skilled workers at given wages.
\end{definition}

\begin{definition}[Canonical model of the skill premium]\label{def:b3:wage-structure:canonical}
\index{canonical model of the skill premium}
The \emph{canonical model} of the \omterm{def:b3:wage-structure:skill-premium}{skill premium} has two kinds of labour, combined
in a constant-elasticity-of-substitution production function, each paid its
marginal product, with the relative demand for skill shifting over time.
\end{definition}

\begin{proposition}[The premium as supply and demand]\label{prop:b3:wage-structure:katz-murphy}
If $Y = [\alpha(A_sS)^\rho + (1 - \alpha)(A_uU)^\rho]^{1/\rho}$ with $\sigma = 1/(1 - \rho)$ and wages equal marginal products, then
$\ln(w_s/w_u) = \tfrac1\sigma[D_t - \ln(S/U)]$, with $D_t = \sigma\ln\tfrac{\alpha}{1-\alpha} + (\sigma - 1)\ln(A_s/A_u)$.
\end{proposition}

\begin{proof}
Differentiating, $\partial Y/\partial S = \alpha A_s^\rho S^{\rho-1}Y^{1-\rho}$ and $\partial Y/\partial U = (1 - \alpha)A_u^\rho U^{\rho-1}Y^{1-\rho}$, so $\ln(w_s/w_u) =
\ln\tfrac{\alpha}{1-\alpha} + \rho\ln(A_s/A_u) + (\rho - 1)\ln(S/U)$. Since $\rho - 1 = -1/\sigma$ and $\rho = (\sigma - 1)/\sigma$, this is $\tfrac1\sigma[\sigma\ln\tfrac{\alpha}{1-\alpha}
+ (\sigma - 1)\ln(A_s/A_u) - \ln(S/U)]$.
\end{proof}

\begin{model}{The race between education and technology}
\assumes Two kinds of workers, a CES production function, and a relative demand
for skill that grows smoothly.
\predicts The premium falls when the supply of skills outpaces demand, and rises
when it lags.
\tested Over 1914--2017 in the United States, a 10~\% rise in the relative supply of
college-educated workers lowers the premium by about 6~\%, implying $\sigma \approx 1.62$, within the
range of 1 to 2.5 found elsewhere; a model with smooth demand trends tracks the premium
closely. % ledger: b3ws:agk.a1
\breaks The fitted demand trend slows after 1992; the returns to education have
become convex, rising most for postgraduate degrees; and wage differences within
education groups, which the model leaves out, grew too. Education differences
explain 75~\% of the rise in wage inequality from 1980 to 2000 but only 38~\% from 2000 to 2017.
\end{model}

\section{Discrimination}\label{sec:b3:wage-structure:discrimination}

\begin{definition}[Taste-based and statistical discrimination]\label{def:b3:wage-structure:discrimination}
\index{taste-based discrimination}\index{statistical discrimination}
\emph{Taste-based discrimination} is unequal treatment arising because employers,
co-workers or customers dislike dealing with members of a group. \emph{Statistical
discrimination} is unequal treatment arising because a group characteristic is
used as a signal of unobserved productivity.
\end{definition}

\begin{proposition}[Competition against taste]\label{prop:b3:wage-structure:becker}
Let all workers have the same productivity, and let prejudiced employers act as if
hiring a worker of group $B$ cost $(1 + d)w_B$. If unprejudiced employers can enter freely
with constant returns, the wage gap $w_A - w_B$ disappears in the long run.
\end{proposition}

\begin{proof}
Suppose $w_B < w_A$. An unprejudiced employer hiring only group $B$ has the same productivity
at a lower wage cost and earns a profit of $w_A - w_B$ per worker above that of an employer
hiring group $A$. With free entry and constant returns, such employers expand, raising
the demand for group $B$ until $w_B = w_A$. With productivity 20 an hour and wages of 20 and 16,
each worker of group $B$ earns the unprejudiced employer 4 an hour: entry continues as long
as any gap remains.
\end{proof}

\begin{model}{Taste-based discrimination}
\assumes Some employers dislike hiring a group and are willing to pay for it.
\predicts Gaps shrink where competition is strong; discriminating firms earn less.
\tested Large firms' racial contact gaps are negatively correlated with their
profitability. % ledger: b3ws:krw
\breaks Gaps persist in correspondence experiments long after any reasonable
time for entry; customers' or co-workers' prejudice, unlike the employer's, is not
competed away.
\end{model}

\begin{definition}[Correspondence (audit) study]\label{def:b3:wage-structure:audit}
\index{correspondence study}\index{audit study}
A \emph{correspondence study} sends fictitious applications, identical except for a
signal of group membership, to real job openings and compares the rates of reply.
\end{definition}

\begin{record}[Names on résumés]
Fictitious résumés sent to help-wanted advertisements in Boston and Chicago with
distinctively white names received 50~\% more callbacks than the same résumés with
distinctively African-American names; a better résumé raised callbacks by 30~\% for
white names and far less for African-American ones. % ledger: b3ws:bm
A later experiment sent more than 83{,}000 applications to 108 of the largest U.S.
employers: distinctively Black names lowered the probability of contact by 2.1
percentage points, with a standard deviation across firms of 1.9 points; firms in the top
fifth accounted for nearly half of the lost contacts. % ledger: b3ws:krw
\end{record}

\begin{definition}[Oaxaca--Blinder decomposition]\label{def:b3:wage-structure:oaxaca}
\index{Oaxaca--Blinder decomposition}
The \emph{Oaxaca--Blinder decomposition} splits the gap in mean log wages between two
groups into a part explained by differences in their measured characteristics and
an unexplained part due to differences in the coefficients.
\end{definition}

\begin{proposition}[The decomposition]\label{prop:b3:wage-structure:oaxaca}
If $y = x'\beta_g + u$ is fitted by OLS with an intercept in each group, then $\bar y_A - \bar y_B = (\bar x_A - \bar x_B)'\beta_A +
\bar x_B'(\beta_A - \beta_B)$. Using $\beta_B$ for the first term instead gives another, generally different, split.
\end{proposition}

\begin{proof}
OLS residuals average zero with an intercept, so $\bar y_g = \bar x_g'\hat\beta_g$. Adding and subtracting $\bar x_B'\hat\beta_A$ gives
the first form; adding and subtracting $\bar x_A'\hat\beta_B$ gives $(\bar x_A - \bar x_B)'\hat\beta_B + \bar x_A'(\hat\beta_A - \hat\beta_B)$. The two differ
unless $\hat\beta_A = \hat\beta_B$: which group's prices value the characteristics is a choice, the
index-number problem.
\end{proof}

\begin{omchart}{Oaxaca--Blinder decomposition of the gap in log 1978 earnings between
white and black men (CPS comparison sample, positive earnings; characteristics:
schooling, age, age squared, married), with white and with black coefficients as
the reference.}{Dehejia--Wahba / LaLonde data (cps\_controls); our regressions}
\begin{tikzpicture}
\begin{axis}[omrecord, xbar stacked, bar width=12pt, width=0.6\linewidth,
  height=0.3\linewidth, xmin=0, xmax=0.36, ymin=-0.6, ymax=1.6, ytick={0,1},
  yticklabels={White coefficients, Black coefficients}, xlabel={Log points},
  x tick label style={/pgf/number format/fixed, /pgf/number format/precision=2},
  scaled x ticks=false, legend columns=2,
  legend style={font=\footnotesize, at={(0.5,-0.32)}, anchor=north}]
  \addplot[fill=omDef!60, draw=omDef, area legend] table[x=explained, y=pos] {figdata/bachelor-3/25-wage-structure/out/oaxaca.csv};
  \addplot[fill=omThm!60, draw=omThm, area legend] table[x=unexplained, y=pos] {figdata/bachelor-3/25-wage-structure/out/oaxaca.csv};
  \legend{Explained, Unexplained}
\end{axis}
\end{tikzpicture}
\end{omchart}

% figdata: oaxaca, oaxaca_stats
Among 11{,}845 white and 977 black men, the gap in mean log earnings is 0.31. Differences in
schooling (12.3 against 11.0 years), age and marriage explain 0.09 of it with white
coefficients and 0.11 with black ones; 0.20 to 0.22 remains. The unexplained part is not a
measure of discrimination: it contains everything the four characteristics miss,
and the characteristics themselves --- schooling among them --- may carry the effects
of earlier discrimination.

\begin{record}[The gender wage gap, 1980--2010]
In the Panel Study of Income Dynamics, women's mean hourly wage rose from 62.1~\% of
men's in 1980 to 79.3~\% in 2010, nearly all of the gain in the 1980s. Blau and Kahn's
decomposition, with male coefficients, shows how the explanation changed. Schooling
and experience accounted for 27~\% of the gap in 1980 and 8~\% in 2010: women had
overtaken men in schooling and nearly closed the gap in experience. Occupation and
industry accounted for 0.117 log points, half the 2010 gap. With all the
characteristics held equal, the ratio rose from 79.4~\% to 91.6~\%, and the
unexplained part fell from 49~\% of the gap to 38~\%. % ledger: b3ws:bk.gap
The occupation share is the same kind of number as the chart's explained part: it
counts differences that may themselves be choices, constraints or the result of
discrimination.
\end{record}

\section{Unions}\label{sec:b3:wage-structure:unions}

\begin{definition}[Union wage premium]\label{def:b3:wage-structure:union}
\index{union wage premium}
The \emph{union wage premium} is the difference between the wage of a unionised
worker and the wage the same worker would earn without a union.
\end{definition}

\begin{measure}[The union premium across the distribution]
Across workers, the measured union wage gap ranges from over 30~\% for lower-paid
workers to 10~\% for higher-paid ones. Following the same workers as they move into and
out of union jobs, and correcting for errors in reported union status, shows that
part of this is a truly larger effect for the lower paid, and part is selection.
The decline of unions among lower-paid workers accounts for a fifth of the rise in the
variance of adult men's wages from 1973 to 1987. % ledger: b3ws:card96
\end{measure}

\section{Migration and local wages}\label{sec:b3:wage-structure:migration}

\begin{definition}[Immigration surplus]\label{def:b3:wage-structure:immigration-surplus}
\index{immigration surplus}
The \emph{immigration surplus} is the net gain to natives from immigration: the
rise in the income of native-owned capital and complementary factors less the
fall in the wages of natives who compete with the immigrants.
\end{definition}

\begin{proposition}[The surplus triangle]\label{prop:b3:wage-structure:surplus}
With labour demand $w = a - bL$, $N$ natives and $M$ immigrants who are perfect substitutes,
all supplying one unit, the \omterm{def:b3:wage-structure:immigration-surplus}{immigration surplus} is $\tfrac12bM^2$; as a share of national income
it is about $\tfrac12 s e m^2$, with $s$ labour's share, $e$ the elasticity of the wage with respect to
employment in absolute value and $m = M/L$ the immigrant share.
\end{proposition}

\begin{proof}
The wage falls from $a - bN$ to $a - b(N + M)$, by $bM$. Natives lose $bMN$ in wages. Output rises by the
area under demand from $N$ to $N + M$; immigrants are paid $(a - b(N + M))M$; the rest of the area, $\tfrac12bM^2$,
and the natives' wage loss $bMN$ go to the owners of the other factors, who are natives.
The net gain is $\tfrac12bM^2$. With $e = bL/w$, $\tfrac12bM^2 = \tfrac12 e(wL)(M/L)^2 = \tfrac12 s e m^2\,Y$.
\end{proof}

\begin{omfigure}
\begin{tikzpicture}[econ, x=0.55cm, y=0.55cm]
  \draw[econ axes] (0,0) -- (11,0) node[right, font=\scriptsize] {$L$};
  \draw[econ axes] (0,0) -- (0,8) node[above, font=\scriptsize] {wage};
  \draw[curve, omDef] (0,7.5) -- (10,1.5) node[right, font=\scriptsize] {labour demand};
  \fill[omDef!15] (0,3.9) rectangle (5,4.5);
  \fill[omThm!45] (5,4.5) -- (6,3.9) -- (5,3.9) -- cycle;
  \draw[guide] (5,0) -- (5,4.5); \draw[guide] (6,0) -- (6,3.9);
  \draw[guide] (0,4.5) -- (5,4.5); \draw[guide] (0,3.9) -- (6,3.9);
  \node[below, font=\scriptsize] at (5,0) {$N$};
  \node[below, font=\scriptsize, xshift=7pt] at (6,0) {$N + M$};
  \node[font=\scriptsize] at (2.5,4.2) {transfer};
  \node[font=\scriptsize, anchor=west] (sp) at (7.2,5.4) {surplus};
  \draw[->, thin] (sp.west) -- (5.35,4.12);
\end{tikzpicture}
\par\smallskip{\small Immigration lowers the wage by $bM$: the rectangle passes from native
workers to native owners of other factors; the triangle is the natives' net gain.}
\end{omfigure}

\begin{debate}[Immigration and local wages]
In 1980 the Mariel boatlift raised Miami's labour force by 7~\%, mostly with
less-skilled workers; the wages of non-Cuban workers over 1979--85 showed virtually no
effect, and less-skilled black workers' unemployment did not rise (\emph{Grade 2}:
difference in differences against comparison cities). % ledger: b3ws:card90
A reappraisal focused on high-school dropouts, at least 60~\% of the arrivals, and found
their wages in Miami fell by 10 to 30~\% (\emph{Grade 2--3}: small cells). % ledger: b3ws:borjas
The difference was traced to a large change, unrelated to the boatlift, in the racial
composition of the small survey samples in Miami (\emph{Grade 2}). % ledger: b3ws:ch
A \omterm{def:b3:difference-in-differences:synthetic-control}{synthetic control} built from cities that matched Miami's trends before the
boatlift finds no significant wage difference for non-Cuban dropouts (\emph{Grade
2}). % ledger: b3ws:py
The reappraisal's author replied that, with the samples adjusted for race, the
low-skill wage in Miami still fell significantly after 1980 against other cities,
and recovered fully by 1990 (\emph{Grade 2--3}: the same small cells). % ledger: b3ws:bor17
Beyond Mariel, national comparisons across cells of education and experience
disagree in the same way: treating immigrants and natives in a cell as perfect
substitutes, a 10~\% rise in a cell's labour supply lowers its wage by 3 to 4~\%
(\emph{Grade 3}); allowing them to be imperfect substitutes, the immigration of
1990--2004 raised the average wage of U.S.-born workers and lowered slightly that of
those without a high-school diploma (\emph{Grade 3}). % ledger: b3ws:bor03, b3ws:op
The Mariel estimates rest on a few dozen workers a year in the Miami survey
samples, and on comparison cities chosen after the fact; the national estimates rest
on assumptions about substitution that the data identify only weakly. Neither says
much about displacement, which can occur without a fall in wages.
\end{debate}

\section{Exercises}\label{sec:b3:wage-structure:exercises}

\begin{exercise}[$\star$]\label{exo:b3:wage-structure:1}
In the opening chart, in which period did the premium rise fastest? Was supply
growth or demand growth more responsible?
\end{exercise}

\begin{exercise}[$\star$]\label{exo:b3:wage-structure:2}
A job carries an extra annual fatality risk of 2 in 10{,}000 and pays 1{,}200 dalens a year more.
Find the implied \omterm{def:b3:cost-benefit-analysis:vsl}{value of a statistical life}.
\end{exercise}

\begin{exercise}[$\star$]\label{exo:b3:wage-structure:3}
Group A has mean schooling 14, group B 12; the return is 0.08 in both. The gap in mean log
wages is 0.25. Decompose it.
\end{exercise}

\begin{exercise}[$\star$]\label{exo:b3:wage-structure:4}
What callback data would contradict \omterm{def:b3:wage-structure:discrimination}{statistical discrimination} by race?
\end{exercise}

\begin{exercise}[$\star\star$]\label{exo:b3:wage-structure:5}
Between two decades the relative supply of graduates rose by 0.30 log points and the
premium fell by 0.05, while relative demand rose by 0.20. Find the implied $\sigma$.
\end{exercise}

\begin{exercise}[$\star\star$]\label{exo:b3:wage-structure:6}
Decompose the white--black earnings gap of the chart variable by variable, with each
group's coefficients. Which characteristic explains most?
\end{exercise}

\begin{exercise}[$\star\star$]\label{exo:b3:wage-structure:7}
A surge of graduates raises the relative supply by 10~\% at given demand. With $\sigma = 1.5$, by how
much does the premium change?
\end{exercise}

\begin{exercise}[$\star\star$]\label{exo:b3:wage-structure:8}
Immigrants make up 10~\% of employment; labour's share is 0.6 and the wage elasticity 0.3. Find the
surplus as a share of national income, and the transfer from native workers.
\end{exercise}

\begin{exercise}[$\star\star$]\label{exo:b3:wage-structure:9}
What evidence would show that the union premium is selection rather than a causal
effect?
\end{exercise}

\begin{exercise}[$\star\star\star$]\label{exo:b3:wage-structure:10}
Derive the canonical-model equation from the CES production function, and show that
\omterm{def:b3:wage-structure:skill-premium}{skill-biased technical change} ($A_s/A_u$ rising) raises the relative demand for skill if and
only if $\sigma > 1$.
\end{exercise}

\begin{exercise}[$\star\star\star$]\label{exo:b3:wage-structure:11}
A prejudiced employer with $d = 0.25$ hires only group A at 20 an hour; group B's wage is 16;
productivity is 20 for both. Compare his profit per worker with an unprejudiced rival's,
and explain why entry closes the gap.
\end{exercise}

\begin{exercise}[$\star\star\star$]\label{exo:b3:wage-structure:12}
In an invented two-city panel the true effect of a shock on wages is zero, but the
treated city's survey sample shifts toward lower-paid workers after the shock.
Compare the difference in differences on raw means with the one within kinds of
worker.
\end{exercise}

\section{Problem: Brenland's Graduates}\label{sec:b3:wage-structure:problem}

\begin{problem}[{Weekend problem --- Brenland's graduates}]\label{pb:b3:wage-structure:graduates}
Brenland's statistical office publishes, for 1990--2025, the log relative supply of
graduates to non-graduates and the log graduate premium (invented series, in the
chapter's \code{brenland_series} table). % figdata: brenland_series, brenland_stats

\medskip\noindent\textbf{Part I --- the premium and supply.}
\begin{enumerate}
\item Report the premium in 1990, 2005 and 2025.
\item Report the log relative supply in the same years.
\item Compute its average annual growth before and after 2005.
\item Describe the premium's path, and say why supply alone cannot explain it.
\item What would a constant premium with rising supply imply about demand?
\end{enumerate}

\medskip\noindent\textbf{Part II --- estimating the model.} Regress the premium on a
time trend and the log relative supply.
\begin{enumerate}[resume]
\item Report the coefficient on supply and its standard error.
\item Derive $\sigma$.
\item Derive the annual growth of relative demand.
\item Compare $\sigma$ with the U.S. estimate.
\item Why does the regression need the time trend?
\item What would bias the coefficient on supply toward zero?
\end{enumerate}

\medskip\noindent\textbf{Part III --- the forecast.} Graduate supply grows 2~\% a year
from 2025 to 2035; demand keeps its trend.
\begin{enumerate}[resume]
\item Compute the relative supply in 2035.
\item Forecast the premium in 2035.
\item Express it as a ratio of wages.
\item How fast would supply need to grow to hold the premium constant?
\item Name a reason the forecast could be too high.
\end{enumerate}

\medskip\noindent\textbf{Part IV --- a gap between two groups.} Group A's mean schooling
and experience are 13 and 18 years, group B's 12 and 20; A's coefficients are 0.08 and
0.02; the gap in mean log wages is 0.15.
\begin{enumerate}[resume]
\item Compute the explained part.
\item Compute the unexplained part.
\item Why does B's longer experience reduce the explained part?
\item What would using B's coefficients require?
\item Name one omitted characteristic that would shrink the unexplained part.
\item Why is the unexplained part not a measure of discrimination?
\item What experiment would measure discrimination directly?
\item Part of the gap may be a graduate premium if A has more graduates. How would you
check?
\item State the result: Brenland's elasticity of substitution between graduates and
non-graduates, and the forecast premium in 2035.
\end{enumerate}
\end{problem}
